\documentclass[11pt]{amsart}

\usepackage[T1]{fontenc}
\usepackage{lmodern}
\usepackage{microtype}
\usepackage{amsmath,amssymb,amsthm}
\usepackage{booktabs}
\usepackage[hidelinks]{hyperref}

\hypersetup{
  pdftitle={A Degree-18 Counterexample to Conjecture 5.23 of Acevedo--Blekherman--Debus--Riener}
}

\newtheorem{theorem}{Theorem}
\newtheorem{lemma}[theorem]{Lemma}

\DeclareMathOperator{\trop}{trop}
\newcommand{\BS}{\mathcal{B}\Sigma}
\newcommand{\EPart}{\Lambda^{\mathrm{ev}}}
\newcommand{\R}{\mathbb{R}}

\title[A degree-18 counterexample]
{A Degree-18 Counterexample to Conjecture~5.23 of
Acevedo--Blekherman--Debus--Riener}

\date{}

\subjclass[2020]{Primary 14T90; Secondary 05E05, 14P10}
\keywords{tropicalization, sums of squares, symmetric functions,
superdominance, partitions}

\begin{document}

\begin{abstract}
Acevedo, Blekherman, Debus, and Riener conjectured that
\(\trop(\BS_{2d}^{*})=T_{2d}^{(2)}\) for every \(d\ge3\).
We disprove their Conjecture~5.23 by giving an explicit vector
\[
y\in\trop(\BS_{18}^{*})\setminus T_{18}^{(2)}.
\]
In exact integer arithmetic, \(y\) satisfies all \(417\) first-order
superdominance inequalities on the even partitions of \(18\) and all
\(1056\) inequalities generated by the degree-\(9\) partial-symmetry
description; by the membership criterion recalled below, this places
\(y\) in the tropicalized sums-of-squares cone. One second-order
superdominance inequality has slack \(-2\).
\end{abstract}

\maketitle

\section{The conjecture and a finite test}

We use the notation of Acevedo, Blekherman, Debus, and
Riener~\cite{ABDR}. Let \(\EPart_{2d}\) be the set of even
partitions of \(2d\). Write the parts of a partition increasingly,
\(
\lambda_{(1)}\le\cdots\le\lambda_{(\ell(\lambda))}
\).
For partitions of the same integer, define superdominance by
\[
\lambda\succeq\mu
\quad\Longleftrightarrow\quad
\sum_{i=1}^{j}\lambda_{(i)}
\le
\sum_{i=1}^{j}\mu_{(i)}
\quad
\text{for }1\le j\le
\min\{\ell(\lambda),\ell(\mu)\}.
\]
If \(\lambda\mu\) denotes multiset union of parts, followed by
sorting, and \(\mu^{\circ2}=\mu\mu\), then
\begin{equation}\label{eq:T2}
T_{2d}^{(2)}=
\left\{
 y\in\R^{\EPart_{2d}}:
 y_{\lambda^1}+y_{\lambda^2}\ge2y_{\mu}
 \text{ whenever }
 \lambda^1\lambda^2\succeq\mu^{\circ2}
\right\}.
\end{equation}
Define also
\[
T_{2d}^{(1)}=
\left\{
 y\in\R^{\EPart_{2d}}:
 y_{\lambda}\ge y_{\mu}
 \text{ whenever }\lambda\succeq\mu
\right\}.
\]
The source proves
\(
T_{2d}^{(1)}\supseteq\trop(\BS_{2d}^{*})\supseteq T_{2d}^{(2)}
\)
and conjectures the second inclusion to be equality for all \(d\ge3\)
\cite[Propositions~5.13 and~5.20 and Conjecture~5.23]{ABDR}; all
numbering here is that of arXiv:2408.04616v1.

We record the finite partial-symmetry test used below. For
\(r,s\ge0\), let \(\mathcal V_{r,s}(d)\) consist of pairs
\(u=(\alpha,\lambda)\), where
\[
\alpha=(\alpha_1,\ldots,\alpha_{r+s})
\in\mathbb Z_{>0}^{r+s},
\]
the first \(r\) entries of \(\alpha\) are odd, the last \(s\) are
even, \(\lambda\) is an even partition, and
\(
|\alpha|+|\lambda|=d
\).
Since \(|\alpha|\equiv r\pmod 2\) and \(|\lambda|\) is even, the parity
constraint forces \(r\equiv d\pmod2\), and \(r+2s\le d\); for \(d=9\)
this leaves \(r\) odd and \(15\) nonempty blocks. All lists of parts
below are sorted into partitions. For \(u=(\alpha,\lambda)\), set
\[
A_u=(2\alpha)\lambda\lambda,
\qquad
B_u=(2\alpha)^{\star}\lambda\lambda
\quad (r+s\ge2),
\]
where \((2\alpha)^{\star}\) is obtained from \((2\alpha)\) by
merging its two largest parts. The restriction \(r+s\ge2\) is not a
choice: for \(r+s=1\) the monomial symmetric function
\(m_{2\alpha}=m_{(2\alpha_1)}\) is the power sum \(p_{2\alpha_1}\),
whose negative support is empty, so single-part blocks contribute no
type-\((1)\) inequality. These are the \(131-119=12\) terms with
\(r+s=1\) counted in the proof of Theorem~\ref{thm:main}. For
\(u=(\alpha,\lambda)\) and \(v=(\beta,\mu)\) in the same block, set
\[
C_{u,v}=(\alpha+\beta)\lambda\mu,
\]
where \(\alpha+\beta\) is componentwise.

\begin{lemma}[Degree-18 membership test]\label{lem:test}
Let \(H\subseteq\R^{\EPart_{18}}\) be the set of vectors \(z\)
satisfying
\begin{equation}\label{eq:H}
 z_{A_u}\ge z_{B_u}
 \quad (u\in\mathcal V_{r,s}(9),\ r+s\ge2),
 \qquad
 z_{A_u}+z_{A_v}\ge2z_{C_{u,v}}
 \quad (u\ne v)
\end{equation}
over all nonempty blocks \(\mathcal V_{r,s}(9)\). Then
\[
\trop(\BS_{18}^{*})=T_{18}^{(1)}\cap H.
\]
Repeated inequalities are harmless.
\end{lemma}

\begin{proof}
We prove \(T_{18}^{(1)}\cap H\subseteq\trop(\BS_{18}^{*})\) first, so
let \(z\in T_{18}^{(1)}\cap H\).
For each \((r,s)\), form the row vector of all terms
\[
n^{(r+s)/2}X_1^{\alpha_1}\cdots
X_{r+s}^{\alpha_{r+s}}p_{\lambda},
\qquad
(\alpha,\lambda)\in\mathcal V_{r,s}(9).
\]
These rows span the degree-\(9\) forms after the partial-symmetry
decomposition. The limiting Gram blocks in
\cite[Section~5.3 and Appendix~A]{ABDR} have, up to positive scalar
factors, entries
\[
 m_{2\alpha}p_{\lambda}^{2}
 \quad\text{and}\quad
 m_{\alpha+\beta}p_{\lambda}p_{\mu}
\]
on the diagonal and off the diagonal, respectively, where \(m_\gamma\)
is a monomial symmetric function. Lemma~5.17 and Corollary~5.18
of~\cite{ABDR} identify the superdominance-maximal indices in the
positive support, the negative support, and the full support of
\(m_{2\alpha}\), respectively \(m_{\alpha+\beta}\). Those corollaries
concern \(m_{2\alpha}\) alone, whereas the indices occurring in the
Gram entries are \((2\alpha)\lambda\lambda\) and
\((\alpha+\beta)\lambda\mu\); the two are matched as follows. The
factor \(p_{\lambda}^{2}\) (resp.\ \(p_{\lambda}p_{\mu}\)) contributes
the same multiset \(\lambda\lambda\) (resp.\ \(\lambda\mu\)) of parts to
every index of the entry, so the merges of Lemma~5.17 take place only
inside \(m_{2\alpha}\) (resp.\ \(m_{\alpha+\beta}\)) and the indices
being compared differ only in that factor. The maximality of
\((2\alpha)^{\star}\) in the negative support of \(m_{2\alpha}\)
therefore transports to \((2\alpha)^{\star}\lambda\lambda=B_u\), and in
the same way the positive diagonal, negative diagonal, and full
off-diagonal maxima are \(A_u\), \(B_u\), and \(C_{u,v}\). We note a
notational divergence: the proof of \cite[Proposition~5.20]{ABDR}
writes the type-\((1)\) inequality coming from a diagonal entry as
\(y_{\nu}\ge y_{\nu^{\star}}\) with \(\nu=\alpha\lambda\), starring the
whole index, whereas \(B_u\) above stars \((2\alpha)\) only.
Because \(z\in T_{18}^{(1)}\), the coordinate maxima are attained at
those same indices, so \eqref{eq:H} verifies the max inequalities
of \cite[Lemma~5.14]{ABDR}. Those
inequalities are quantified over all pairs of rows, whereas
\eqref{eq:H} imposes an off-diagonal condition only for \(u\ne v\)
lying in the same block; for rows in different blocks the entry of the
limiting Gram matrix vanishes, so the associated index set is empty and
the condition is vacuous under the convention
\(\max\emptyset=-\infty\) of \cite[Lemma~5.14]{ABDR}. The
full-dimensional closure argument in the proof of
\cite[Proposition~5.19]{ABDR} then places \(z\) in
\(\trop(\BS_{18}^{*})\).

For the reverse inclusion, \(\trop(\BS_{18}^{*})\subseteq T_{18}^{(1)}\)
by \cite[Proposition~5.13]{ABDR}, and every point of
\(\trop(\BS_{18}^{*})\) satisfies the max inequalities
of \cite[Lemma~5.14]{ABDR}. On \(T_{18}^{(1)}\) the maxima in those
inequalities are attained at \(A_u\), \(B_u\), and \(C_{u,v}\), as
above, so there they reduce to \eqref{eq:H}. Hence
\(\trop(\BS_{18}^{*})\subseteq T_{18}^{(1)}\cap H\).
\end{proof}

\section{The counterexample}

Define \(y\in\R^{\EPart_{18}}\) by the following table; exponential
notation denotes repeated parts.
{\small
\begin{equation}\label{eq:witness}
\renewcommand{\arraystretch}{1.04}
\begin{array}{
  c@{\,}r@{\qquad}
  c@{\,}r@{\qquad}
  c@{\,}r
}
\toprule
\lambda & y_\lambda &
\lambda & y_\lambda &
\lambda & y_\lambda \\
\midrule
(18)       & 0  & (10,4,2^2) & 24 & (6^2,2^3)   & 36  \\
(16,2)     & 9  & (10,2^4)   & 44 & (6,4^3)     & 20  \\
(14,4)     & 6  & (8^2,2)    & 12 & (6,4^2,2^2) & 30  \\
(14,2^2)   & 20 & (8,6,4)    & 10 & (6,4,2^4)   & 48  \\
(12,6)     & 3  & (8,6,2^2)  & 24 & (6,2^6)     & 68  \\
(12,4,2)   & 14 & (8,4^2,2)  & 21 & (4^4,2)     & 30  \\
(12,2^3)   & 32 & (8,4,2^3)  & 36 & (4^3,2^3)   & 40  \\
(10,8)     & 0  & (8,2^5)    & 56 & (4^2,2^5)   & 60  \\
(10,6,2)   & 12 & (6^3)      & 10 & (4,2^7)     & 80  \\
(10,4^2)   & 12 & (6^2,4,2)  & 20 & (2^9)       & 100 \\
\bottomrule
\end{array}
\end{equation}
}

\begin{theorem}\label{thm:main}
The vector \(y\) in \eqref{eq:witness} satisfies
\[
y\in\trop(\BS_{18}^{*})\setminus T_{18}^{(2)}.
\]
Consequently, Conjecture~5.23 of~\cite{ABDR} is false and
\[
\trop(\BS_{18}^{*})\supsetneq T_{18}^{(2)}.
\]
\end{theorem}

\begin{proof}
For \(d=9\), there are \(15\) nonempty blocks
\(\mathcal V_{r,s}(9)\), containing \(131\) terms in total.
There are \(417\) nontrivial first-order superdominance
inequalities on \(\EPart_{18}\), and exact substitution gives minimum
slack \(0\) over all of them. The vector \(y\) therefore belongs to
\(T_{18}^{(1)}\). The system \eqref{eq:H} consists of \(119\) diagonal
and \(937\) off-diagonal inequalities, hence \(1056\) generated
inequalities; exact substitution again gives minimum slack \(0\), so
\(y\in H\). By Lemma~\ref{lem:test},
\[
y\in\trop(\BS_{18}^{*}).
\]

Now take
\[
\lambda^1=(10,8),
\qquad
\lambda^2=(4,4,4,2,2,2),
\qquad
\mu=(8,4,4,2).
\]
The increasingly ordered prefix sums of
\(\lambda^1\lambda^2\) and \(\mu^{\circ2}\) are
\[
(2,4,6,10,14,18,26,36)
\quad\text{and}\quad
(2,4,8,12,16,20,28,36),
\]
respectively. Hence
\(
\lambda^1\lambda^2\succeq\mu^{\circ2}
\), so \eqref{eq:T2} requires
\[
y_{(10,8)}+y_{(4,4,4,2,2,2)}
\ge2y_{(8,4,4,2)}.
\]
Its slack at \eqref{eq:witness} is
\[
0+40-2\cdot21=-2.
\]
Therefore \(y\notin T_{18}^{(2)}\), as claimed.
\end{proof}

\paragraph{Exact verification.}
The program reproduced in the appendix uses
only the Python standard library and no floating-point arithmetic. It
generates all partitions, blocks,
and inequalities from the definitions above, checks the counts
\(15,131,119,937,1056\), verifies all \(417\) first-order
superdominance inequalities and all partial-symmetry inequalities with
minimum slack \(0\), and verifies the second-order violation with
slack \(-2\).
It records two further facts about the finite test of
Lemma~\ref{lem:test}: each of the \(119\) diagonal inequalities has
\(A_u\succeq B_u\) and is therefore already a first-order inequality,
and each of the \(937\) off-diagonal ones has
\(A_uA_v\succeq C_{u,v}^{\circ2}\) and is therefore itself an instance
of~\eqref{eq:T2}. Finally, of the \(7025\) inequalities
of~\eqref{eq:T2} on \(\EPart_{18}\)---one for each unordered pair
\(\{\lambda^1,\lambda^2\}\), repetition allowed, together with each
\(\mu\) satisfying \(\lambda^1\lambda^2\succeq\mu^{\circ2}\)---exactly
one is violated by \(y\), namely
\(y_{(10,8)}+y_{(4,4,4,2,2,2)}\ge2y_{(8,4,4,2)}\).

\raggedbottom

\appendix

\section{The verification program}\label{app:code}

The following self-contained program performs the checks described
above. Comments give the expected output, one
\texttt{print} statement per line.

{\small
\begin{verbatim}
from itertools import combinations
from itertools import combinations_with_replacement as cwr

def parts(n, mx=None):            # partitions, parts decreasing
    mx = n if mx is None else mx
    if n == 0:
        yield ()
        return
    for k in range(min(n, mx), 0, -1):
        for t in parts(n - k, k):
            yield (k,) + t

def evens(m):                     # even partitions of m
    return [] if m % 2 else [tuple(2*p for p in t)
                             for t in parts(m//2)]

def key(ls):
    return tuple(sorted(ls, reverse=True))

def sd(a, b):                     # a superdominates b
    A, B, sa, sb = sorted(a), sorted(b), 0, 0
    for i in range(min(len(A), len(B))):
        sa, sb = sa + A[i], sb + B[i]
        if sa > sb:
            return False
    return True

def alphas(r, s, m):              # ordered: r odd, then s even
    if r:
        for a in range(1, m + 1, 2):
            for t in alphas(r - 1, s, m - a):
                yield (a,) + t
    elif s:
        for a in range(2, m + 1, 2):
            for t in alphas(0, s - 1, m - a):
                yield (a,) + t
    else:
        yield ()

Y = {(18,):0,(16,2):9,(14,4):6,(14,2,2):20,(12,6):3,
 (12,4,2):14,(12,2,2,2):32,(10,8):0,(10,6,2):12,(10,4,4):12,
 (10,4,2,2):24,(10,2,2,2,2):44,(8,8,2):12,(8,6,4):10,
 (8,6,2,2):24,(8,4,4,2):21,(8,4,2,2,2):36,(8,2,2,2,2,2):56,
 (6,6,6):10,(6,6,4,2):20,(6,6,2,2,2):36,(6,4,4,4):20,
 (6,4,4,2,2):30,(6,4,2,2,2,2):48,(6,2,2,2,2,2,2):68,
 (4,4,4,4,2):30,(4,4,4,2,2,2):40,(4,4,2,2,2,2,2):60,
 (4,2,2,2,2,2,2,2):80,(2,2,2,2,2,2,2,2,2):100}

d, L = 9, evens(18)
assert len(L) == 30 and sorted(L) == sorted(Y)
one = [(a, b) for a in L for b in L if a != b and sd(a, b)]
print(len(one), min(Y[a] - Y[b] for a, b in one))    # 417 0

blk = {}
for r in range(d + 1):
    for s in range(d//2 + 1):
        V = [(al, la) for al in alphas(r, s, d)
             for la in evens(d - sum(al))] if r + s else []
        if V:
            blk[(r, s)] = V
print(len(blk), sum(len(V) for V in blk.values()))   # 15 131

A = lambda u: key([2*a for a in u[0]] + 2*list(u[1]))
def Bs(u):
    t = sorted([2*a for a in u[0]], reverse=True)
    return key([t[0] + t[1]] + t[2:] + 2*list(u[1]))
C = lambda u, v: key([a + b for a, b in zip(u[0], v[0])]
                     + list(u[1]) + list(v[1]))

dg = [(A(u), Bs(u)) for (r, s), V in blk.items()
      if r + s >= 2 for u in V]
of = [(A(u), A(v), C(u, v)) for V in blk.values()
      for u, v in combinations(V, 2)]
print(len(dg), len(of), len(dg) + len(of))       # 119 937 1056
print(min(Y[a] - Y[b] for a, b in dg),
      min(Y[a] + Y[b] - 2*Y[c] for a, b, c in of))   # 0 0
print(all(sd(a, b) for a, b in dg),
      all(sd(a + b, c + c) for a, b, c in of))       # True True

two = [(a, b, m) for a, b in cwr(L, 2) for m in L
       if sd(a + b, m + m)]
bad = [t for t in two if Y[t[0]] + Y[t[1]] - 2*Y[t[2]] < 0]
print(len(two), bad)
# 7025 [((10, 8), (4, 4, 4, 2, 2, 2), (8, 4, 4, 2))]
print(Y[(10,8)] + Y[(4,4,4,2,2,2)] - 2*Y[(8,4,4,2)])      # -2
\end{verbatim}
}

\begin{thebibliography}{1}

\bibitem{ABDR}
J.~Acevedo, G.~Blekherman, S.~Debus, and C.~Riener,
\emph{Symmetric nonnegative functions, the tropical Vandermonde
cell and superdominance of power sums},
Collect. Math. \textbf{77} (2026), 555--600,
\href{https://doi.org/10.1007/s13348-025-00481-z}
{doi:10.1007/s13348-025-00481-z};
\href{https://arxiv.org/abs/2408.04616}{arXiv:2408.04616}.

\end{thebibliography}

\end{document}